% Lean: EconLean/Macro/DP/{Blackwell,SequenceProblem,Bellman,Concavity,BrockMirman,Examples}.lean
\chapter{Dynamic Programming}
\label{chap:dp}

Source: Stokey--Lucas (SL), \emph{Recursive Methods in Economic Dynamics}, ch.~3--4.
A deterministic dynamic program has a state space $X$, a feasibility correspondence
$\Gamma : X \twoheadrightarrow X$, a one-period return $F(x, y)$ and a discount factor $\beta$.

\section{Blackwell's sufficient conditions}
\label{sec:blackwell}

We work in the space $C_b(X)$ of bounded continuous real functions on a topological space $X$
with the sup distance (SL's $B(X)$ is the case of the discrete topology).

\begin{theorem}[Blackwell; SL Thm~3.3]
\label{thm:blackwell}
\lean{Econ.blackwell}
\leanok
Let $0 \le \beta < 1$ and $T : C_b(X) \to C_b(X)$ satisfy
(monotonicity) $f \le g \Rightarrow Tf \le Tg$ pointwise, and
(discounting) $T(f + a) \le Tf + \beta a$ pointwise for every constant $a \ge 0$.
Then $T$ is a contraction of modulus $\beta$.
\end{theorem}
\begin{proof}
\leanok
Let $a = d(f, g) = \|f - g\|_\infty \ge 0$. For any pair $f, g$ at distance $a$, we have
$f \le g + a$ pointwise, so by monotonicity and then discounting
$Tf \le T(g + a) \le Tg + \beta a$. Applying this to $(f, g)$ and to $(g, f)$ (same distance) gives
$|Tf(x) - Tg(x)| \le \beta a$ for every $x$, hence $d(Tf, Tg) \le \beta\, d(f, g)$.
\end{proof}

\begin{corollary}[Contraction mapping theorem for Blackwell operators; SL Thm~3.2]
\label{cor:blackwell_fixedPoint}
\lean{Econ.blackwell_fixedPoint}
\leanok
\uses{thm:blackwell}
Under the hypotheses of Theorem~\ref{thm:blackwell}, $T$ has a unique fixed point $v$, and for every
$v_0$ and $n$, $d(T^n v_0, v) \le d(v_0, T v_0)\, \beta^n / (1 - \beta)$.
\end{corollary}
\begin{proof}
\leanok
$C_b(X)$ is complete, so Banach's fixed point theorem (Mathlib's \texttt{ContractingWith.fixedPoint}
and its a priori estimate) applies to the contraction of Theorem~\ref{thm:blackwell}. Uniqueness is
uniqueness of the contraction's fixed point.
\end{proof}

\section{The sequence problem}
\label{sec:sequence}

This section makes no topological assumptions, so it covers returns that are unbounded below, such as
log utility. To avoid Lean's junk values for divergent series, payoffs are handled through partial sums.

\begin{definition}[Plans, payoffs, optimality]
\label{def:sequence_problem}
\lean{Econ.IsFeasiblePlan, Econ.partialPayoff, Econ.IsOptimalValue, Econ.IsOptimalPlan}
\leanok
A \emph{feasible plan} from $x_0$ is a sequence $(x_t)$ with $x_0$ given and $x_{t+1} \in \Gamma(x_t)$.
Its \emph{partial payoff} is $u_n(x) = \sum_{t < n} \beta^t F(x_t, x_{t+1})$.
A number $w$ is the \emph{optimal value} at $x_0$ (attained) if some feasible plan has $u_n \to w$ and
every feasible plan has $\limsup_n u_n \le w$ (i.e.\ for each $\varepsilon > 0$, eventually
$u_n \le w + \varepsilon$). A feasible plan $x$ is \emph{optimal} if for every feasible $x'$ and
$\varepsilon > 0$, eventually $u_n(x') \le u_n(x) + \varepsilon$.
\end{definition}

\begin{lemma}[Basic properties of optimal values]
\label{lem:optimal_value_props}
\lean{Econ.IsOptimalValue.unique, Econ.IsOptimalValue.isOptimalPlan, Econ.IsOptimalValue.tendsto_of_isOptimalPlan}
\leanok
\uses{def:sequence_problem}
(1) The optimal value is unique. (2) A feasible plan whose payoffs converge to the optimal value is an
optimal plan. (3) If $w$ is the optimal value, every optimal plan has $u_n \to w$.
\end{lemma}
\begin{proof}
\leanok
(1) If $w, w'$ are both optimal values, take the plan with $u_n \to w$; since eventually
$u_n \le w' + \varepsilon$, we get $w \le w' + \varepsilon$ for all $\varepsilon > 0$, so $w \le w'$; by
symmetry $w = w'$.
(2) Given feasible $x'$ and $\varepsilon$: eventually $u_n(x') \le w + \varepsilon/2$, and since
$u_n(x) \to w$, eventually $w - \varepsilon/2 < u_n(x)$; add.
(3) Let $x^s$ be the plan attaining $w$. Eventually: $u_n(x) \le w + \varepsilon/2$ (as $w$ is the
optimal value); $u_n(x^s) \le u_n(x) + \varepsilon/3$ (as $x$ is optimal); and
$w - \varepsilon/3 < u_n(x^s)$. Together $|u_n(x) - w| < \varepsilon$.
\end{proof}

\begin{definition}[Policy plan]
\label{def:planOf}
\lean{Econ.planOf, Econ.planOf_succ, Econ.isFeasiblePlan_planOf}
\leanok
\uses{def:sequence_problem}
For $g : X \to X$, the plan generated from $x_0$ is $x_t = g^t(x_0)$, so $x_{t+1} = g(x_t)$. If
$g(x) \in \Gamma(x)$ for all $x$, it is feasible.
\end{definition}

\begin{lemma}[Telescoping identity]
\label{lem:telescope}
\lean{Econ.telescope}
\leanok
\uses{def:sequence_problem}
For any $v : X \to \R$, plan $x$ and $n$,
\[
  v(x_0) = u_n(x) + \beta^n v(x_n) + \sum_{t < n} \beta^t \big(v(x_t) - F(x_t, x_{t+1}) - \beta v(x_{t+1})\big).
\]
\end{lemma}
\begin{proof}
\leanok
Induction on $n$. For $n = 0$ both sides are $v(x_0)$. The step adds the date-$n$ terms
$\beta^n F(x_n, x_{n+1})$ and $\beta^n(v(x_n) - F(x_n, x_{n+1}) - \beta v(x_{n+1}))$ and replaces
$\beta^n v(x_n)$ by $\beta^{n+1} v(x_{n+1})$; the result is a ring identity.
\end{proof}

\begin{lemma}[Bellman inequality and equality along a plan]
\label{lem:partialPayoff_bellman}
\lean{Econ.partialPayoff_add_le, Econ.partialPayoff_add_eq}
\leanok
\uses{lem:telescope}
If $\beta \ge 0$ and $F(x_t, x_{t+1}) + \beta v(x_{t+1}) \le v(x_t)$ for all $t$, then
$u_n + \beta^n v(x_n) \le v(x_0)$. With equality at every $t$, $u_n + \beta^n v(x_n) = v(x_0)$.
\end{lemma}
\begin{proof}
\leanok
In Lemma~\ref{lem:telescope}, each gap term is $\beta^t$ times a nonnegative number (resp.\ zero).
\end{proof}

\begin{lemma}[Invariant sets]
\label{lem:feasible_mem}
\lean{Econ.IsFeasiblePlan.mem}
\leanok
\uses{def:sequence_problem}
If $\Gamma(x) \subseteq S$ for all $x \in S$ and $x_0 \in S$, every feasible plan from $x_0$ stays in $S$.
\end{lemma}
\begin{proof}
\leanok
Induction on $t$.
\end{proof}

\begin{theorem}[Verification theorem; SL Thms~4.3--4.4]
\label{thm:verification}
\lean{Econ.verification}
\leanok
\uses{lem:partialPayoff_bellman, lem:feasible_mem, def:planOf}
Let $\beta \ge 0$, $S$ invariant under $\Gamma$, and $v : X \to \R$, $g : X \to X$ satisfy on $S$:
(i) $F(x, y) + \beta v(y) \le v(x)$ for $y \in \Gamma(x)$;
(ii) $g(x) \in \Gamma(x)$ and $F(x, g(x)) + \beta v(g(x)) = v(x)$;
(iii) for every feasible plan from $x_0 \in S$ and $\varepsilon > 0$, eventually
$-\varepsilon \le \beta^n v(x_n)$;
(iv) along the plan generated by $g$, eventually $\beta^n v(x_n) \le \varepsilon$.
Then for $x_0 \in S$, $v(x_0)$ is the optimal value and the plan generated by $g$ attains it.
\end{theorem}
\begin{proof}
\leanok
\emph{Upper bound.} For a feasible plan $x$, it stays in $S$ (Lemma~\ref{lem:feasible_mem}), so by (i)
and Lemma~\ref{lem:partialPayoff_bellman}, $u_n \le v(x_0) - \beta^n v(x_n)$, and by (iii) this is
eventually $\le v(x_0) + \varepsilon$.
\emph{The policy plan.} Let $x^g$ be generated by $g$. By induction it stays in $S$ (using (ii) and
invariance), so it is feasible, and by (ii) and Lemma~\ref{lem:partialPayoff_bellman},
$u_n(x^g) + \beta^n v(x^g_n) = v(x_0)$ for all $n$.
\emph{Convergence.} Given $\varepsilon$, eventually $u_n(x^g) \le v(x_0) + \varepsilon/2$ (upper bound)
and $\beta^n v(x^g_n) \le \varepsilon/2$ (by (iv)), so
$u_n(x^g) = v(x_0) - \beta^n v(x^g_n) \ge v(x_0) - \varepsilon/2$. Hence $u_n(x^g) \to v(x_0)$, which
with the upper bound shows $v(x_0)$ is the optimal value.
\end{proof}

\begin{lemma}[A geometric-type recursion tends to zero]
\label{lem:tendsto_zero_rec}
\lean{Econ.tendsto_zero_of_le_mul_add_pow}
\leanok
If $s_n \ge 0$, $s_{n+1} \le r s_n + d \beta^n$ with $r < 1$, $d \ge 0$ and $0 \le \beta < 1$, then
$s_n \to 0$.
\end{lemma}
\begin{proof}
\leanok
Let $q = \max(r, \beta, \tfrac12) \in (0, 1)$. By induction, $s_n \le (s_0 + n\, d/q)\, q^n$: the step uses
$r s_n \le q (s_0 + n d/q) q^n$ and $d \beta^n \le d q^n$, and
$(s_0 + (n+1) d/q) q^{n+1} = q (s_0 + n d/q) q^n + d q^n$. The bound tends to $0$ because $q^n \to 0$
and $n q^n \to 0$; squeeze.
\end{proof}

\begin{theorem}[Verification for returns unbounded below]
\label{thm:verification_lowerBound}
\lean{Econ.verification_of_lowerBound}
\leanok
\uses{lem:partialPayoff_bellman, lem:feasible_mem, lem:tendsto_zero_rec, def:planOf}
Let $0 \le \beta < 1$ and replace (iii) in Theorem~\ref{thm:verification} by a lower bound
$a + b F(x, y) \le v(x)$ for $x \in S$, $y \in \Gamma(x)$, with $b \ge 1$. Then the same conclusion holds.
\end{theorem}
\begin{proof}
\leanok
Only the upper bound changes. Fix a feasible plan $x$ and let $s_n = \max(u_n - v(x_0), 0) \ge 0$.
At date $n$: Lemma~\ref{lem:partialPayoff_bellman} gives $u_n + \beta^n v(x_n) \le v(x_0)$; the lower
bound times $\beta^n \ge 0$ gives $\beta^n a + \beta^n b F(x_n, x_{n+1}) \le \beta^n v(x_n)$; and
$u_{n+1} = u_n + \beta^n F(x_n, x_{n+1})$. Combining,
$b (u_{n+1} - v(x_0)) \le (b - 1)(u_n - v(x_0)) - \beta^n a \le (b - 1) s_n + |a| \beta^n$.
Dividing by $b > 0$ and noting the right side is $\ge 0$,
$s_{n+1} \le \frac{b-1}{b} s_n + \frac{|a|}{b} \beta^n$. Since $\frac{b-1}{b} < 1$,
Lemma~\ref{lem:tendsto_zero_rec} gives $s_n \to 0$, so eventually $u_n - v(x_0) \le s_n < \varepsilon$.
The policy plan is handled exactly as in Theorem~\ref{thm:verification}.
\end{proof}

\begin{theorem}[Euler equations and transversality are sufficient; SL Thm~4.15]
\label{thm:euler_transversality}
\lean{Econ.euler_transversality_along, Econ.euler_transversality}
\leanok
\uses{def:sequence_problem}
Let states lie in an ordered real vector space $E$, with feasible next states nonnegative. Let
$F_x(x, y), F_y(x, y)$ be linear functionals with the supergradient inequality on the graph of $\Gamma$,
\[
  F(x', y') \le F(x, y) + F_x(x, y)(x' - x) + F_y(x, y)(y' - y),
\]
and $F_x(x, y) z \ge 0$ for $z \ge 0$. If a feasible plan $x^*$ satisfies the Euler equations
$F_y(x^*_t, x^*_{t+1}) + \beta F_x(x^*_{t+1}, x^*_{t+2}) = 0$ and the transversality condition
$\beta^t F_x(x^*_t, x^*_{t+1}) x^*_t \to 0$, and $\beta \ge 0$, then $x^*$ is optimal.
The Lean proof is \texttt{euler\_transversality\_along}, which only asks for the supergradient inequality and
$F_x \ge 0$ at the points $(x^*_t, x^*_{t+1})$ of the candidate plan; \texttt{euler\_transversality} is the
global form above, derived from it.
\end{theorem}
\begin{proof}
\leanok
Let $x$ be feasible, write $D^x_t = F_x(x^*_t, x^*_{t+1})$, $D^y_t = F_y(x^*_t, x^*_{t+1})$ and
$z_t = x^*_t - x_t$, so $z_0 = 0$.
\emph{Step 1 (concavity).} The supergradient inequality at $(x^*_t, x^*_{t+1})$ evaluated at
$(x_t, x_{t+1})$ gives $D^x_t z_t + D^y_t z_{t+1} \le F(x^*_t, x^*_{t+1}) - F(x_t, x_{t+1})$.
\emph{Step 2 (Euler).} By induction on $n$,
$\sum_{t<n} \beta^t (D^x_t z_t + D^y_t z_{t+1}) = -\beta^n D^x_n z_n$: the step uses the Euler equation
applied to $z_{n+1}$, $D^y_n z_{n+1} = -\beta D^x_{n+1} z_{n+1}$.
\emph{Step 3.} Multiplying Step 1 by $\beta^t \ge 0$ and summing,
$-\beta^n D^x_n z_n \le u_n(x^*) - u_n(x)$. For $n \ge 1$, $x_n \ge 0$ (it is a feasible next state), so
$D^x_n x_n \ge 0$ and $-\beta^n D^x_n z_n \ge -\beta^n D^x_n x^*_n$. Hence
$u_n(x) \le u_n(x^*) + \beta^n D^x_n x^*_n$.
\emph{Step 4.} By transversality, eventually $\beta^n D^x_n x^*_n < \varepsilon$.
\end{proof}

\section{The Bellman equation with bounded returns}
\label{sec:bellman}

\begin{definition}[Dynamic program; SL Assumptions 4.3--4.4]
\label{def:DynProg}
\lean{Econ.DynProg, Econ.DynProg.objective, Econ.DynProg.bound}
\leanok
\uses{def:contCorr}
A \emph{dynamic program} on a topological space $X$ consists of $\Gamma$ with nonempty compact values that
is a continuous correspondence, a continuous $F : X \times X \to \R$ bounded on the graph of $\Gamma$ by
some $M$, and $0 < \beta < 1$. Its Bellman objective is $F(x, y) + \beta v(y)$.
\end{definition}

\begin{lemma}[The objective is continuous and the maximum is attained]
\label{lem:dp_objective}
\lean{Econ.DynProg.continuous_objective, Econ.DynProg.abs_F_le, Econ.DynProg.isGreatest_maxValue_objective}
\leanok
\uses{def:DynProg, lem:isGreatest_maxValue}
For $v \in C_b(X)$, $(x, y) \mapsto F(x, y) + \beta v(y)$ is continuous, and its maximum over $\Gamma(x)$
is attained.
\end{lemma}
\begin{proof}
\leanok
Sum and products of continuous functions; Lemma~\ref{lem:isGreatest_maxValue}.
\end{proof}

\begin{definition}[Bellman operator]
\label{def:bellmanOp}
\lean{Econ.DynProg.bellmanOp, Econ.DynProg.bellmanOp_apply, Econ.DynProg.isGreatest_bellmanOp}
\leanok
\uses{lem:dp_objective, thm:continuous_maxValue}
$(Tv)(x) = \max_{y \in \Gamma(x)} [F(x, y) + \beta v(y)]$. It maps $C_b(X)$ to itself: $Tv$ is continuous
by Berge (Theorem~\ref{thm:continuous_maxValue}), and $|Tv(x)| \le M + \beta \|v\|$ since at a maximizer
$y$, $|F(x, y) + \beta v(y)| \le |F(x, y)| + \beta |v(y)|$.
\end{definition}

\begin{theorem}[$T$ is a contraction; SL Thm~4.6]
\label{thm:bellman_contracting}
\lean{Econ.DynProg.contracting}
\leanok
\uses{def:bellmanOp, thm:blackwell}
$T$ is a contraction of modulus $\beta$.
\end{theorem}
\begin{proof}
\leanok
Check Blackwell's conditions (Theorem~\ref{thm:blackwell}).
Monotone: if $f \le g$, let $y$ attain $Tf(x)$; then
$Tf(x) = F(x, y) + \beta f(y) \le F(x, y) + \beta g(y) \le Tg(x)$.
Discounting: let $y$ attain $T(f + a)(x)$; then
$T(f + a)(x) = F(x, y) + \beta f(y) + \beta a \le Tf(x) + \beta a$.
\end{proof}

\begin{theorem}[Existence and uniqueness of the value function; SL Thm~4.6]
\label{thm:bellman_equation}
\lean{Econ.DynProg.value, Econ.DynProg.bellmanOp_value, Econ.DynProg.bellman_equation, Econ.DynProg.value_unique, Econ.DynProg.iterate_converges}
\leanok
\uses{thm:bellman_contracting}
Let $V$ be the fixed point of $T$. Then $V(x) = \max_{y \in \Gamma(x)} [F(x, y) + \beta V(y)]$ (attained),
$V$ is the only bounded continuous solution of this equation, and for any $v_0$,
$T^n v_0 \to V$ with $d(T^n v_0, V) \le d(v_0, T v_0) \beta^n / (1 - \beta)$.
\end{theorem}
\begin{proof}
\leanok
$V$ exists by Banach's theorem; $TV = V$ and Definition~\ref{def:bellmanOp} give the equation. If $w$
solves the equation, then $Tw(x)$ and $w(x)$ are both the greatest element of the same set, so $Tw = w$
and $w = V$ by uniqueness of fixed points. The convergence and error bound are Banach's a priori estimate.
\end{proof}

\begin{definition}[Policy correspondence]
\label{def:policy}
\lean{Econ.DynProg.policy, Econ.DynProg.mem_policy_iff, Econ.DynProg.policy_subset}
\leanok
\uses{thm:bellman_equation, def:maxValue}
$G(x) = \operatorname{argmax}_{y \in \Gamma(x)} [F(x, y) + \beta V(y)]$. Then $y \in G(x)$ iff
$y \in \Gamma(x)$ and $F(x, y) + \beta V(y) = V(x)$.
\end{definition}
\begin{proof}
\leanok
($\Rightarrow$) $F(x, y) + \beta V(y) \le V(x)$ by the Bellman equation, and $\ge$ the value at a
maximizer $y_0$, which is $V(x)$. ($\Leftarrow$) for any $y' \in \Gamma(x)$,
$F(x, y') + \beta V(y') \le V(x) = F(x, y) + \beta V(y)$.
\end{proof}

\begin{theorem}[Policy correspondence is nonempty, compact and UHC; SL Thm~4.6]
\label{thm:policy_uhc}
\lean{Econ.DynProg.policy_nonempty_compact_uhc, Econ.DynProg.policy_nonempty, Econ.DynProg.continuous_policy_of_singleton}
\leanok
\uses{def:policy, lem:argmax_props, thm:uhc_argmax, cor:continuous_argmax_singleton}
$G$ has nonempty compact values and is UHC; if $G(x) = \{g(x)\}$ for all $x$, then $g$ is continuous.
\end{theorem}
\begin{proof}
\leanok
Berge's theorem (Lemma~\ref{lem:argmax_props}, Theorem~\ref{thm:uhc_argmax},
Corollary~\ref{cor:continuous_argmax_singleton}) applied to the objective with $v = V$.
\end{proof}

\begin{definition}[Payoff of a plan]
\label{def:dp_payoff}
\lean{Econ.DynProg.payoff, Econ.DynProg.summable_payoff, Econ.DynProg.tendsto_partialPayoff}
\leanok
\uses{def:DynProg, def:sequence_problem}
The payoff of a plan is $\sum_t \beta^t F(x_t, x_{t+1})$. For a feasible plan the series converges
absolutely, being dominated by $M \beta^t$, so its partial sums $u_n$ converge to it.
\end{definition}

\begin{lemma}[Discounted values vanish]
\label{lem:tendsto_discounted_value}
\lean{Econ.DynProg.tendsto_discounted_value, Econ.DynProg.bellman_le}
\leanok
\uses{thm:bellman_equation}
For any sequence $(x_n)$, $\beta^n V(x_n) \to 0$; and $F(x, y) + \beta V(y) \le V(x)$ for $y \in \Gamma(x)$.
\end{lemma}
\begin{proof}
\leanok
$|\beta^n V(x_n)| \le \beta^n \|V\| \to 0$. The inequality is the Bellman equation.
\end{proof}

\begin{theorem}[No plan beats $V$; plans following $G$ attain $V$; SL Thms~4.3--4.4]
\label{thm:payoff_value}
\lean{Econ.DynProg.payoff_le_value, Econ.DynProg.payoff_eq_value_of_policy}
\leanok
\uses{lem:partialPayoff_bellman, lem:tendsto_discounted_value, def:dp_payoff, def:policy}
For a feasible plan from $x_0$, $\text{payoff} \le V(x_0)$; if moreover $x_{t+1} \in G(x_t)$ for all $t$,
then $\text{payoff} = V(x_0)$.
\end{theorem}
\begin{proof}
\leanok
By Lemma~\ref{lem:partialPayoff_bellman}, $u_n \le V(x_0) - \beta^n V(x_n)$, whose right side tends to
$V(x_0)$ (Lemma~\ref{lem:tendsto_discounted_value}); pass to the limit. Along a policy plan, equality
holds at each date (Definition~\ref{def:policy}), so $u_n = V(x_0) - \beta^n V(x_n) \to V(x_0)$; limits
are unique.
\end{proof}

\begin{theorem}[Optimal plans follow the policy; SL Thm~4.5]
\label{thm:policy_of_optimal}
\lean{Econ.DynProg.policy_of_payoff_eq_value}
\leanok
\uses{lem:telescope, lem:tendsto_discounted_value, def:policy}
If a feasible plan has payoff $V(x_0)$, then $x_{t+1} \in G(x_t)$ for every $t$.
\end{theorem}
\begin{proof}
\leanok
Let $\mathrm{gap}_t = V(x_t) - F(x_t, x_{t+1}) - \beta V(x_{t+1}) \ge 0$ and
$G_n = \sum_{t<n} \beta^t \mathrm{gap}_t$. By Lemma~\ref{lem:telescope},
$G_n = V(x_0) - u_n - \beta^n V(x_n)$, which tends to $V(x_0) - \text{payoff} - 0 = 0$. $G_n$ is
nondecreasing and tends to $0$, so $G_n \le 0$; also $G_n \ge 0$. Then
$\beta^t \mathrm{gap}_t = G_{t+1} - G_t$ is both $\le 0$ and $\ge 0$, so $\mathrm{gap}_t = 0$ (as
$\beta^t > 0$), which is the membership criterion of Definition~\ref{def:policy}.
\end{proof}

\begin{theorem}[Principle of optimality; SL Thms~4.2--4.5]
\label{thm:principle_of_optimality}
\lean{Econ.DynProg.policyPlan, Econ.DynProg.policyPlan_spec, Econ.DynProg.principle_of_optimality}
\leanok
\uses{thm:payoff_value, thm:policy_of_optimal, thm:policy_uhc}
$V(x_0)$ is the greatest payoff over feasible plans from $x_0$, and a feasible plan attains it iff
$x_{t+1} \in G(x_t)$ for all $t$.
\end{theorem}
\begin{proof}
\leanok
Choose a selection $g(x) \in G(x)$ (nonempty by Theorem~\ref{thm:policy_uhc}); the plan it generates is
feasible and follows $G$, so it attains $V(x_0)$ by Theorem~\ref{thm:payoff_value}, which also gives the
upper bound. The equivalence is Theorems~\ref{thm:payoff_value} and~\ref{thm:policy_of_optimal}.
\end{proof}

\begin{corollary}[Agreement with the sequence-problem notion]
\label{cor:isOptimalValue_value}
\lean{Econ.DynProg.isOptimalValue_value}
\leanok
\uses{thm:verification, lem:tendsto_discounted_value, thm:policy_uhc}
$V(x_0)$ is the optimal value at $x_0$ in the sense of Definition~\ref{def:sequence_problem}.
\end{corollary}
\begin{proof}
\leanok
Apply Theorem~\ref{thm:verification} with $S = X$ and a selection $g$ from $G$. Conditions (iii) and (iv)
hold because $\beta^n V(x_n) \to 0$ along every sequence (Lemma~\ref{lem:tendsto_discounted_value}).
\end{proof}

\begin{theorem}[Monotone value function; SL Thm~4.7]
\label{thm:value_monotone}
\lean{Econ.DynProg.value_monotone, Econ.DynProg.value_strictMono}
\leanok
\uses{thm:bellman_equation, lem:tendsto_discounted_value}
Let $X$ be preordered and $\Gamma$ monotone ($x \le x' \Rightarrow \Gamma(x) \subseteq \Gamma(x')$). If
$F(\cdot, y)$ is nondecreasing (strictly increasing), then $V$ is nondecreasing (strictly increasing).
\end{theorem}
\begin{proof}
\leanok
Let $x \le x'$ (resp.\ $x < x'$) and $y$ attain $V(x)$. Then $y \in \Gamma(x')$, so
$V(x) = F(x, y) + \beta V(y) \le F(x', y) + \beta V(y) \le V(x')$ (with strict first inequality in the
strict case).
\end{proof}

\section{Concavity, unique policies, and Benveniste--Scheinkman}
\label{sec:dp_concavity}

The state space is a convex set $S$ in a real normed space $E$, and the dynamic program lives on $S$.

\begin{definition}[Convex technology, strictly concave returns; SL Assumptions 4.7--4.8]
\label{def:dp_concavity}
\lean{Econ.ConvexGraph, Econ.StrictlyConcaveReturn, Econ.StrictlyConcaveFn, Econ.ConcaveFn}
\leanok
$\Gamma$ has \emph{convex graph} if $y \in \Gamma(x)$, $y' \in \Gamma(x')$, $a, b > 0$, $a + b = 1$ imply
$a y + b y' \in \Gamma(a x + b x')$. $F$ is \emph{strictly concave on the graph} if for such data
$a F(x, y) + b F(x', y') \le F(ax + bx', ay + by')$, strictly when $x \neq x'$. A function $f$ on $S$
is concave (strictly concave) if $a f(x) + b f(x') \le f(ax + bx')$ (strictly when $x \neq x'$).
\end{definition}

\begin{lemma}[Agreement with Mathlib's concavity]
\label{lem:concaveFn_concaveOn}
\lean{Econ.extendS, Econ.ConcaveFn.concaveOn}
\leanok
\uses{def:dp_concavity}
Extend $f : S \to \R$ by $0$ outside $S$. If $f$ is concave in the sense above, the extension is
\texttt{ConcaveOn} $S$.
\end{lemma}
\begin{proof}
\leanok
For $a, b \ge 0$ with $a + b = 1$: if both are positive this is the hypothesis; if $a = 0$ then $b = 1$
and both sides equal $f(y)$; symmetrically if $b = 0$.
\end{proof}

\begin{lemma}[$T$ maps concave functions to strictly concave ones; SL Thm~4.8, key step]
\label{lem:strictlyConcave_bellmanOp}
\lean{Econ.DynProg.strictlyConcave_bellmanOp}
\leanok
\uses{def:dp_concavity, def:bellmanOp}
Under convex technology and strictly concave returns, if $f$ is concave then $Tf$ is strictly concave.
\end{lemma}
\begin{proof}
\leanok
Let $z = ax + bx'$, and let $y, y'$ attain $Tf(x), Tf(x')$. Then $w = ay + by' \in \Gamma(z)$ by convex
technology, so $Tf(z) \ge F(z, w) + \beta f(w)$. Now
\[
  a\,Tf(x) + b\,Tf(x') = [a F(x, y) + b F(x', y')] + \beta [a f(y) + b f(y')],
\]
and the first bracket is $\le F(z, w)$ (strictly if $x \neq x'$), the second is $\le f(w)$ by
concavity of $f$, and $\beta > 0$.
\end{proof}

\begin{theorem}[Strictly concave value function; SL Thm~4.8]
\label{thm:value_strictlyConcave}
\lean{Econ.DynProg.value_strictlyConcave}
\leanok
\uses{lem:strictlyConcave_bellmanOp, thm:bellman_equation}
Under convex technology and strictly concave returns, $V$ is strictly concave.
\end{theorem}
\begin{proof}
\leanok
By induction every iterate $T^n 0$ is concave (the zero function is concave;
Lemma~\ref{lem:strictlyConcave_bellmanOp} for the step). $T^n 0 \to V$ uniformly
(Theorem~\ref{thm:bellman_equation}), and for fixed $x, x', z, a, b$ the map
$f \mapsto f(z) - (a f(x) + b f(x'))$ is continuous, so its nonnegativity passes to the limit: $V$ is
concave. Then $V = TV$ is strictly concave by Lemma~\ref{lem:strictlyConcave_bellmanOp}.
\end{proof}

\begin{theorem}[Unique, continuous policy function; SL Thm~4.8]
\label{thm:exists_continuous_policy}
\lean{Econ.DynProg.exists_continuous_policy}
\leanok
\uses{thm:value_strictlyConcave, def:policy, thm:policy_uhc}
Under convex technology and strictly concave returns there is a continuous $g : S \to S$ with
$G(x) = \{g(x)\}$ for all $x$.
\end{theorem}
\begin{proof}
\leanok
Uniqueness: suppose $y \neq y'$ are both in $G(x)$, so $F(x, y) + \beta V(y) = V(x) = F(x, y') + \beta V(y')$.
The midpoint $w = \tfrac12 y + \tfrac12 y'$ is feasible at $x = \tfrac12 x + \tfrac12 x$ (convex
technology), $F(x, w) \ge \tfrac12 F(x, y) + \tfrac12 F(x, y')$, and $V(w) > \tfrac12 V(y) + \tfrac12 V(y')$
by strict concavity of $V$ (Theorem~\ref{thm:value_strictlyConcave}) as $y \neq y'$. So
$F(x, w) + \beta V(w) > V(x)$, contradicting the Bellman inequality. Choose $g(x) \in G(x)$ (nonempty);
then $G(x) = \{g(x)\}$ and $g$ is continuous by Theorem~\ref{thm:policy_uhc}.
\end{proof}

\begin{lemma}[Concave function with a differentiable minorant; SL Thm~4.10, one dimension]
\label{lem:hasDerivAt_concave_minorant}
\lean{Econ.hasDerivAt_of_concaveOn_of_le}
\leanok
Let $V$ be concave on a neighbourhood $D$ of $x_0 \in \R$, $W$ differentiable at $x_0$ with derivative
$w$, $W \le V$ near $x_0$ and $W(x_0) = V(x_0)$. Then $V$ is differentiable at $x_0$ with $V'(x_0) = w$.
\end{lemma}
\begin{proof}
\leanok
Work with difference quotients. Pick $\delta > 0$ with $(x_0 - \delta, x_0 + \delta) \subseteq D$ and
$W \le V$ there. Since $W(x_0) = V(x_0)$:
for $x > x_0$, $\frac{W(x) - W(x_0)}{x - x_0} \le \frac{V(x) - V(x_0)}{x - x_0}$; for $x < x_0$ the
inequality reverses. By concavity (slopes decrease), for $x < x_0 < x''$,
$\frac{V(x_0) - V(x)}{x_0 - x} \ge \frac{V(x'') - V(x_0)}{x'' - x_0}$.
\emph{Left limit.} For $x < x_0$: the slope of $V$ on $[x, x_0]$ is $\le$ the slope of $W$, which tends
to $w$; and it is $\ge$ every right slope of $V$, hence $\ge$ every right slope of $W$, whose limit is
$w$. Squeeze: the left slope of $V$ tends to $w$.
\emph{Right limit.} Symmetric: right slopes of $V$ are $\ge$ those of $W$ ($\to w$) and $\le$ every left
slope of $V$, hence $\le$ every left slope of $W$ ($\to w$). The two one-sided limits agree, so $V$ is
differentiable at $x_0$ with derivative $w$.
\end{proof}

\begin{theorem}[Benveniste--Scheinkman; SL Thm~4.11, one-dimensional state]
\label{thm:benveniste_scheinkman}
\lean{Econ.DynProg.hasDerivAt_value}
\leanok
\uses{lem:hasDerivAt_concave_minorant, thm:value_strictlyConcave, lem:concaveFn_concaveOn, def:policy}
Let $S \subseteq \R$ be convex, with convex technology and strictly concave returns. Let $x_0$ be interior to
$S$, $y_0 \in G(x_0)$, and suppose $y_0 \in \Gamma(x)$ for all $x \in S$ near $x_0$. If $F(\cdot, y_0)$ is
differentiable at $x_0$ with derivative $F_x$, then $V$ is differentiable at $x_0$ and $V'(x_0) = F_x$.
\end{theorem}
\begin{proof}
\leanok
$V$ is concave on $S$ (Theorem~\ref{thm:value_strictlyConcave} and Lemma~\ref{lem:concaveFn_concaveOn}),
and $S$ is a neighbourhood of $x_0$. Define the minorant $W(x) = F(x, y_0) + \beta V(y_0)$: keep
tomorrow's choice fixed and vary today's state. Near $x_0$, $y_0 \in \Gamma(x)$, so $W(x) \le V(x)$ by the
Bellman inequality; and $W(x_0) = V(x_0)$ because $y_0 \in G(x_0)$. $W$ has derivative $F_x$ at $x_0$.
Apply Lemma~\ref{lem:hasDerivAt_concave_minorant}.
\end{proof}

\begin{remark}[Non-vacuity of the dynamic-programming hypotheses]
\label{rem:dp_example}
\lean{Econ.DPExample.S, Econ.DPExample.F, Econ.DPExample.P, Econ.DPExample.convexGraph, Econ.DPExample.strictlyConcaveReturn}
\leanok
\uses{def:DynProg, def:dp_concavity, thm:value_monotone, thm:exists_continuous_policy, thm:benveniste_scheinkman}
\texttt{Macro/DP/Examples.lean} exhibits $S = [0, 1]$, $\Gamma(x) = S$, $F(x, y) = -(x - 2)^2 - y^2$,
$\beta = 1/2$. It satisfies SL Assumptions 4.3--4.4 ($|F| \le 5$), the monotonicity hypotheses of
Theorem~\ref{thm:value_monotone}, convex technology, and strict concavity, using the identity
$F(ax + bx', ay + by') - [a F(x, y) + b F(x', y')] = ab[(x - x')^2 + (y - y')^2]$. Hence
Theorems~\ref{thm:exists_continuous_policy} and~\ref{thm:benveniste_scheinkman} apply, the latter giving
$V'(1/2) = 3$.
\end{remark}

\section{The Brock--Mirman model}
\label{sec:brock_mirman}

The planner solves $\max \sum_t \beta^t \log c_t$ with $c_t = k_t^\alpha - k_{t+1}$, $0 < \alpha, \beta < 1$,
$k_0 > 0$.

\begin{definition}[Brock--Mirman primitives]
\label{def:brock_mirman}
\lean{Econ.BrockMirman.Γ, Econ.BrockMirman.F, Econ.BrockMirman.A, Econ.BrockMirman.B, Econ.BrockMirman.V, Econ.BrockMirman.g}
\leanok
\uses{def:sequence_problem}
Feasibility $\Gamma(k) = \{k' \mid k > 0,\ 0 < k' < k^\alpha\}$; return $F(k, k') = \log(k^\alpha - k')$;
closed form $V(k) = A + B \log k$ with $B = \frac{\alpha}{1 - \alpha\beta}$ and
$A = \frac{1}{1-\beta}\big[\log(1 - \alpha\beta) + \frac{\alpha\beta}{1-\alpha\beta}\log(\alpha\beta)\big]$;
policy $g(k) = \alpha\beta k^\alpha$.
\end{definition}

\begin{lemma}[Elementary bounds]
\label{lem:bm_ab}
\lean{Econ.BrockMirman.ab_pos, Econ.BrockMirman.ab_lt_one, Econ.BrockMirman.Γ_subset}
\leanok
\uses{def:brock_mirman}
$0 < \alpha\beta < 1$, and every feasible $k'$ is positive.
\end{lemma}

\begin{lemma}[Bellman equation at the policy]
\label{lem:bm_bellman_eq}
\lean{Econ.BrockMirman.bellman_eq_at_policy}
\leanok
\uses{lem:bm_ab}
For $k > 0$, $F(k, g(k)) + \beta V(g(k)) = V(k)$.
\end{lemma}
\begin{proof}
\leanok
$k^\alpha - \alpha\beta k^\alpha = (1 - \alpha\beta) k^\alpha$. Expand
$\log((1-\alpha\beta)k^\alpha) = \log(1-\alpha\beta) + \alpha\log k$ and
$\log(\alpha\beta k^\alpha) = \log(\alpha\beta) + \alpha \log k$; the identity is then algebra in
$\log k$, $\log(1-\alpha\beta)$, $\log(\alpha\beta)$ using the definitions of $A$ and $B$.
\end{proof}

\begin{lemma}[Bellman inequality]
\label{lem:bm_bellman_le}
\lean{Econ.BrockMirman.bellman_le}
\leanok
\uses{lem:bm_bellman_eq}
For every $k' \in \Gamma(k)$, $F(k, k') + \beta V(k') \le V(k)$.
\end{lemma}
\begin{proof}
\leanok
Rewrite $V(k)$ by Lemma~\ref{lem:bm_bellman_eq}; we must show
$\log c - \log c^* + \beta B (\log k' - \log k^*) \le 0$ where $c = k^\alpha - k'$,
$c^* = (1 - \alpha\beta) k^\alpha$, $k^* = \alpha\beta k^\alpha$. By $\log z \le z - 1$,
$\log c - \log c^* \le c/c^* - 1$ and $\log k' - \log k^* \le k'/k^* - 1$. With weight
$\beta B = \frac{\alpha\beta}{1-\alpha\beta} \ge 0$ on the second, the combined linear bound
$\frac{c}{c^*} - 1 + \beta B(\frac{k'}{k^*} - 1)$ is identically $0$ (field algebra).
\end{proof}

\begin{lemma}[Lower bound by the current return]
\label{lem:bm_lowerBound}
\lean{Econ.BrockMirman.lowerBound}
\leanok
\uses{lem:bm_ab}
For $k' \in \Gamma(k)$, $A + \frac{1}{1-\alpha\beta} F(k, k') \le V(k)$.
\end{lemma}
\begin{proof}
\leanok
$0 < k^\alpha - k' \le k^\alpha$ gives $\log(k^\alpha - k') \le \alpha \log k$. Multiply by
$\frac{1}{1 - \alpha\beta} > 0$ and note $B \log k = \frac{1}{1-\alpha\beta}(\alpha\log k)$.
\end{proof}

\begin{lemma}[The policy plan stays in a log-bounded region]
\label{lem:bm_policy_bounds}
\lean{Econ.BrockMirman.policy_plan_bounds}
\leanok
\uses{lem:bm_ab, def:planOf}
Along $k_{t+1} = g(k_t)$ from $k_0 > 0$: $k_t > 0$ and $|\log k_t| \le |\log k_0| + \frac{|\log \alpha\beta|}{1-\alpha}$.
\end{lemma}
\begin{proof}
\leanok
Induction. $k_{t+1} = \alpha\beta k_t^\alpha > 0$ and
$|\log k_{t+1}| = |\log\alpha\beta + \alpha \log k_t| \le |\log\alpha\beta| + \alpha(|\log k_0| + C)$
with $C = |\log\alpha\beta|/(1-\alpha)$; since $|\log\alpha\beta| = (1-\alpha) C$ and $\alpha \le 1$, this is
$\le |\log k_0| + C$.
\end{proof}

\begin{theorem}[Brock--Mirman closed form, value route]
\label{thm:brockMirman_value}
\lean{Econ.BrockMirman.brockMirman_value}
\leanok
\uses{thm:verification_lowerBound, lem:bm_bellman_eq, lem:bm_bellman_le, lem:bm_lowerBound, lem:bm_policy_bounds, lem:optimal_value_props}
For $0 < \alpha, \beta < 1$ and $k_0 > 0$, $A + B\log k_0$ is the optimal value of the sequence problem,
the plan $k_{t+1} = \alpha\beta k_t^\alpha$ is an optimal plan, and its partial payoffs converge to
$A + B \log k_0$.
\end{theorem}
\begin{proof}
\leanok
Apply Theorem~\ref{thm:verification_lowerBound} with $S = (0, \infty)$ (invariant by
Lemma~\ref{lem:bm_ab}), $v = V$, the policy $g$ (feasible: $0 < \alpha\beta k^\alpha < k^\alpha$), the
Bellman inequality and equality (Lemmas~\ref{lem:bm_bellman_le}, \ref{lem:bm_bellman_eq}) and the lower
bound with $a = A$, $b = \frac{1}{1-\alpha\beta} \ge 1$ (Lemma~\ref{lem:bm_lowerBound}). For the policy
condition, Lemma~\ref{lem:bm_policy_bounds} bounds $V(k_n) \le C := |A| + B(|\log k_0| + |\log\alpha\beta|/(1-\alpha))$,
so $\beta^n V(k_n) \le \beta^n C \to 0$. Optimality of the plan follows from Lemma~\ref{lem:optimal_value_props}(2).
\end{proof}

\begin{definition}[Marginal returns]
\label{def:bm_marginal}
\lean{Econ.BrockMirman.Fk, Econ.BrockMirman.Fk'}
\leanok
\uses{def:brock_mirman}
$F_k(k, k') = \alpha k^{\alpha-1}/(k^\alpha - k')$ and $F_{k'}(k, k') = -1/(k^\alpha - k')$, as linear maps on $\R$.
\end{definition}

\begin{lemma}[Supergradient inequality]
\label{lem:bm_supergradient}
\lean{Econ.BrockMirman.supergradient}
\leanok
\uses{def:bm_marginal}
For $y \in \Gamma(k)$ and $y' \in \Gamma(k')$,
$F(k', y') \le F(k, y) + F_k(k, y)(k' - k) + F_{k'}(k, y)(y' - y)$.
\end{lemma}
\begin{proof}
\leanok
Two concavity facts. (Bernoulli) $k'^\alpha \le k^\alpha + \alpha k^{\alpha - 1}(k' - k)$: write
$k' = k(1 + s)$ with $s = k'/k - 1 \ge -1$ and use $(1 + s)^\alpha \le 1 + \alpha s$ for $0 \le \alpha \le 1$.
(Log) $\log z \le \log c + (z - c)/c$ with $c = k^\alpha - y$, $z = k'^\alpha - y'$, from
$\log(z/c) \le z/c - 1$. Then $(z - c)/c = (k'^\alpha - k^\alpha - (y' - y))/c \le
[\alpha k^{\alpha-1}(k' - k) - (y' - y)]/c$, which is the right-hand side.
\end{proof}

\begin{theorem}[Euler equation and transversality along the policy]
\label{thm:bm_euler_tvc}
\lean{Econ.BrockMirman.brockMirman_euler_equation, Econ.BrockMirman.brockMirman_transversality}
\leanok
\uses{lem:bm_policy_bounds}
Along $k_{t+1} = \alpha\beta k_t^\alpha$ with $c_t = k_t^\alpha - k_{t+1}$:
$\frac{1}{c_t} = \frac{\beta \alpha k_{t+1}^{\alpha - 1}}{c_{t+1}}$, and
$\beta^t \frac{\alpha k_t^{\alpha-1}}{c_t} k_t \to 0$.
\end{theorem}
\begin{proof}
\leanok
Along the plan $c_t = (1 - \alpha\beta) k_t^\alpha$ and $k_t > 0$ (Lemma~\ref{lem:bm_policy_bounds}).
Euler: substitute $c_t$, $c_{t+1}$, $k_{t+1}^{\alpha-1} = k_{t+1}^\alpha / k_{t+1}$ and
$k_{t+1} = \alpha\beta k_t^\alpha$; both sides reduce to $\frac{1}{(1-\alpha\beta)k_t^\alpha}$.
Transversality: the term equals $\beta^t \frac{\alpha}{1-\alpha\beta}$, which tends to $0$.
\end{proof}

\begin{theorem}[Brock--Mirman, Euler route]
\label{thm:brockMirman_optimal_of_euler}
\lean{Econ.BrockMirman.brockMirman_optimal_of_euler}
\leanok
\uses{thm:euler_transversality, lem:bm_supergradient, thm:bm_euler_tvc}
The plan $k_{t+1} = \alpha\beta k_t^\alpha$ is optimal, independently of Theorem~\ref{thm:brockMirman_value}.
\end{theorem}
\begin{proof}
\leanok
Apply Theorem~\ref{thm:euler_transversality} with $E = \R$ and the marginal returns of
Definition~\ref{def:bm_marginal}. The supergradient inequality is Lemma~\ref{lem:bm_supergradient};
feasible next states are positive; $F_k \ge 0$ since $k^{\alpha-1} > 0$ and $k^\alpha - k' > 0$. The plan is
feasible. The Euler equation $F_{k'}(k_t, k_{t+1}) + \beta F_k(k_{t+1}, k_{t+2}) = 0$ and transversality
$\beta^t F_k(k_t, k_{t+1}) k_t = \beta^t \alpha/(1-\alpha\beta) \to 0$ are the computations of
Theorem~\ref{thm:bm_euler_tvc}.
\end{proof}

% Lean: EconLean/Macro/DP/Monotone.lean
\section{Monotone policies and monotone convergence}
\label{sec:dp_monotone}

Sources: SL \S4.2, \S6.1; Topkis (1978).

\begin{theorem}[Monotone policy (Topkis)]
\label{thm:policy_monotone}
\lean{Econ.DynProg.policy_monotone}
\leanok
\uses{def:policy, lem:tendsto_discounted_value}
Let $X$ be linearly ordered. Suppose feasible sets are ascending (if $x < x'$, $y \in \Gamma(x)$, $y' \in \Gamma(x')$
and $y' < y$, then $y' \in \Gamma(x)$ and $y \in \Gamma(x')$), and $F$ has strictly increasing differences (for
$x < x'$, $y' < y$, all four pairs feasible, $F(x', y') - F(x, y') < F(x', y) - F(x, y)$). If $y$ is optimal at $x$,
$y'$ is optimal at $x' > x$, then $y \le y'$.
\end{theorem}
\begin{proof}
\leanok
Suppose $y' < y$. By ascending sets, $y' \in \Gamma(x)$ and $y \in \Gamma(x')$. Optimality of $y$ at $x$ against $y'$
(\texttt{bellman\_le}) gives $F(x, y') + \beta V(y') \le F(x, y) + \beta V(y)$; optimality of $y'$ at $x'$ against
$y$ gives $F(x', y) + \beta V(y) \le F(x', y') + \beta V(y')$. Adding, $F(x', y) - F(x, y) \le F(x', y') - F(x, y')$,
contradicting strictly increasing differences.
\end{proof}

\begin{lemma}[Orbits of monotone maps are monotone]
\label{lem:orbit_monotone}
\lean{Econ.orbit_monotone_or_antitone}
\leanok
\uses{def:planOf}
If $g$ is monotone on a linear order, the orbit $x_{n+1} = g(x_n)$ is monotone if $x_0 \le g(x_0)$ and antitone
otherwise.
\end{lemma}
\begin{proof}
\leanok
If $x_0 \le x_1$, induct: $x_n \le x_{n+1}$ gives $x_{n+1} = g(x_n) \le g(x_{n+1}) = x_{n+2}$. Otherwise
$x_1 \le x_0$ and the same induction gives $x_{n+2} \le x_{n+1}$.
\end{proof}

\begin{theorem}[Monotone convergence to a steady state]
\label{thm:orbit_tendsto_fixedPoint}
\lean{Econ.orbit_tendsto_fixedPoint}
\leanok
\uses{lem:orbit_monotone}
If $g : [a, b] \to [a, b]$ is monotone and continuous, every orbit converges to some $k$ with $g(k) = k$.
\end{theorem}
\begin{proof}
\leanok
By Lemma~\ref{lem:orbit_monotone} the orbit is monotone or antitone, and it is bounded by $a$ and $b$, so it converges
to its supremum or infimum $L$ (\texttt{tendsto\_atTop\_ciSup} / \texttt{ciInf}). $[a, b]$ is closed, so $L \in [a, b]$.
The shifted sequence $x_{n+1} = g(x_n)$ tends to $g(L)$ by continuity and to $L$ as a tail of $x$; limits are unique.
\end{proof}

% Lean: EconLean/Macro/DP/MarkovChain.lean
\section{Finite Markov chains}
\label{sec:markov}

Sources: SL ch.~11 (Thms 11.1, 11.4); Ljungqvist--Sargent ch.~2.

\begin{definition}[Stochastic matrices and invariant distributions]
\label{def:markov}
\lean{Econ.IsStochastic, Econ.Distributions, Econ.IsInvariant, Econ.l1Dist, Econ.vecMul_apply'}
\leanok
$Q$ is stochastic if $Q(z, z') \ge 0$ and each row sums to $1$. Distributions are $\mu \ge 0$ with $\sum_z \mu(z) = 1$;
$\mu$ moves to $(\mu Q)(z') = \sum_z \mu(z) Q(z, z')$. $\pi$ is invariant if it is a distribution with $\pi Q = \pi$.
$\|\mu - \nu\|_1 = \sum_z |\mu(z) - \nu(z)|$ (twice the total variation distance).
\end{definition}

\begin{lemma}[Distributions]
\label{lem:distributions}
\lean{Econ.IsStochastic.vecMul_mem, Econ.IsStochastic.pow_vecMul_mem, Econ.isCompact_distributions, Econ.l1Dist_le_two, Econ.l1Dist_self, Econ.l1Dist_eq_zero}
\leanok
\uses{def:markov}
(i) $\mu Q$ and $\mu Q^n$ are distributions when $\mu$ is. (ii) The set of distributions is compact. (iii) For
distributions, $\|\mu - \nu\|_1 \le 2$, $\|\mu - \mu\|_1 = 0$, and $\|\mu - \nu\|_1 = 0$ implies $\mu = \nu$.
\end{lemma}
\begin{proof}
\leanok
(i) Entries are sums of nonnegative products; swapping sums, $\sum_{z'} \sum_z \mu(z) Q(z, z') = \sum_z \mu(z) = 1$.
Powers by induction. (ii) It is a closed subset (an intersection of half-spaces and a hyperplane) of the compact cube
$[0, 1]^Z$; each $\mu(z) \le \sum \mu = 1$. (iii) $|\mu(z) - \nu(z)| \le \mu(z) + \nu(z)$, which sums to $2$; a sum of
nonnegative terms is zero only if each term is.
\end{proof}

\begin{theorem}[Existence of an invariant distribution]
\label{thm:exists_invariant}
\lean{Econ.exists_invariant}
\leanok
\uses{lem:distributions}
Every stochastic matrix on a finite nonempty set has an invariant distribution.
\end{theorem}
\begin{proof}
\leanok
Krylov--Bogolyubov. Start from a point mass $\mu_0$, let $a_t = \mu_0 Q^t$ and the Ces\`aro averages
$\nu_n = \frac{1}{n+1}\sum_{t \le n} a_t$, which are distributions. By compactness some subsequence
$\nu_{\varphi(k)} \to \pi$, a distribution. Shifting the sum, $\nu_n Q - \nu_n = \frac{1}{n+1}(a_{n+1} - a_0)$, whose
entries are bounded by $1/(n+1) \to 0$. By continuity of $\mu \mapsto \mu Q - \mu$, $\nu_{\varphi(k)} Q - \nu_{\varphi(k)}
\to \pi Q - \pi$, and also $\to 0$; so $\pi Q = \pi$.
\end{proof}

\begin{lemma}[Doeblin contraction]
\label{lem:l1_contract}
\lean{Econ.l1Dist_contract}
\leanok
\uses{def:markov}
If $Q(z, z') \ge \varepsilon$ for all $z, z'$, then for distributions $\mu, \nu$,
$\|\mu Q - \nu Q\|_1 \le (1 - |Z|\varepsilon)\|\mu - \nu\|_1$.
\end{lemma}
\begin{proof}
\leanok
Let $d = \mu - \nu$, so $\sum_z d(z) = 0$. Subtracting the common floor,
$(\mu Q - \nu Q)(z') = \sum_z d(z)(Q(z, z') - \varepsilon)$. Then
$\sum_{z'} |\cdot| \le \sum_{z'}\sum_z |d(z)|(Q(z, z') - \varepsilon) = \sum_z |d(z)| (1 - |Z|\varepsilon)$, using that
rows sum to $1$.
\end{proof}

\begin{theorem}[Doeblin's theorem, finite case]
\label{thm:doeblin}
\lean{Econ.doeblin}
\leanok
\uses{thm:exists_invariant, lem:l1_contract, lem:distributions}
If $Q$ is stochastic and $Q(z, z') \ge \varepsilon > 0$ for all $z, z'$, there is a unique invariant distribution $\pi$,
and $\|\mu Q^n - \pi\|_1 \le 2(1 - |Z|\varepsilon)^n$ for every distribution $\mu$. (SL Thm 11.4 uses an $N$-step
condition; this is the one-step case.)
\end{theorem}
\begin{proof}
\leanok
Take $\pi$ from Theorem~\ref{thm:exists_invariant} and let $r = 1 - |Z|\varepsilon$. A row sums to $1$ and has $|Z|$
entries $\ge \varepsilon$, so $r \ge 0$; $r < 1$ since $|Z| > 0$. Using $\pi Q^n = \pi$ and
Lemma~\ref{lem:l1_contract}, induction gives $\|\mu Q^n - \pi\|_1 \le r^n \|\mu - \pi\|_1 \le 2 r^n$. Uniqueness: another
invariant $\pi'$ has $\pi' Q^n = \pi'$, so $\|\pi' - \pi\|_1 \le 2 r^n \to 0$, hence $\|\pi' - \pi\|_1 = 0$ and $\pi' = \pi$.
\end{proof}

% Lean: EconLean/Macro/DP/Stochastic.lean, StochasticExamples.lean
\section{Stochastic dynamic programming with finite shocks}
\label{sec:stoch_dp}

Source: SL ch.~9 (Thm 9.6), with a finite shock space $Z$ (SL allows a compact metric space with a Feller kernel). The
sequence-problem side (SL Thms 9.2--9.5) is not yet formalized, so $V$ is characterized as the unique bounded continuous
solution of the Bellman equation.

\begin{definition}[Stochastic dynamic program]
\label{def:stoch_dp}
\lean{Econ.StochDynProg, Econ.StochDynProg.expect, Econ.StochDynProg.objective, Econ.StochDynProg.bound, Econ.StochDynProg.abs_F_le}
\leanok
\uses{def:contCorr}
States are $(x, z)$ with $z$ in a finite discrete $Z$. $Q$ is a stochastic matrix, $0 < \beta < 1$, $\Gamma(x, z)$ is
nonempty, compact and continuous, and $F((x, z), y)$ is continuous and bounded by $M$ on the graph. For bounded continuous
$v$, $E[v \mid z](y) = \sum_{z'} Q(z, z') v(y, z')$ and the objective is $F((x, z), y) + \beta E[v \mid z](y)$.
\end{definition}

\begin{lemma}[Objective is continuous and bounded]
\label{lem:stoch_objective}
\lean{Econ.StochDynProg.continuous_objective, Econ.StochDynProg.abs_expect_le, Econ.StochDynProg.isGreatest_objective}
\leanok
\uses{def:stoch_dp, lem:isGreatest_maxValue}
The objective is jointly continuous, $|E[v \mid z](y)| \le \|v\|$, and its maximum over $\Gamma(x, z)$ is attained.
\end{lemma}
\begin{proof}
\leanok
$z \mapsto Q(z, z')$ is continuous because $Z$ is discrete, so the objective is a finite sum of continuous terms.
$|\sum_{z'} Q(z, z') v(y, z')| \le \sum_{z'} Q(z, z') \|v\| = \|v\|$. Attainment is Lemma~\ref{lem:isGreatest_maxValue}.
\end{proof}

\begin{theorem}[Stochastic Bellman operator is a contraction]
\label{thm:stoch_contracting}
\lean{Econ.StochDynProg.bellmanOp, Econ.StochDynProg.isGreatest_bellmanOp, Econ.StochDynProg.contracting}
\leanok
\uses{lem:stoch_objective, thm:continuous_maxValue, thm:blackwell}
$(Tv)(x, z) = \max_{y \in \Gamma(x,z)} [F((x,z), y) + \beta E[v \mid z](y)]$ maps bounded continuous functions to bounded
continuous functions, with $\|Tv\| \le M + \beta\|v\|$, and is a contraction of modulus $\beta$.
\end{theorem}
\begin{proof}
\leanok
$Tv$ is continuous by Theorem~\ref{thm:continuous_maxValue}, and bounded since at a maximizer
$|F + \beta E| \le M + \beta \|v\|$. Contraction by Blackwell (Theorem~\ref{thm:blackwell}). Monotonicity: if $f \le g$,
at a maximizer $y$ for $f$, $\sum Q f \le \sum Q g$ (nonnegative weights), so $(Tf)(s) \le$ objective of $g$ at $y \le (Tg)(s)$.
Discounting: rows sum to $1$, so $E[f + a \mid z] = E[f \mid z] + a$; at a maximizer for $f + a$ the value is
$\le (Tf)(s) + \beta a$.
\end{proof}

\begin{theorem}[Stochastic Bellman equation, SL Thm 9.6]
\label{thm:stoch_bellman}
\lean{Econ.StochDynProg.value, Econ.StochDynProg.bellman_equation, Econ.StochDynProg.value_unique, Econ.StochDynProg.iterate_converges}
\leanok
\uses{thm:stoch_contracting}
The fixed point $V$ of $T$ solves the Bellman equation with the max attained, is the only bounded continuous solution,
and value iteration converges: $T^n v_0 \to V$ with $d(T^n v_0, V) \le d(v_0, T v_0)\, \beta^n/(1 - \beta)$.
\end{theorem}
\begin{proof}
\leanok
Banach's fixed point theorem (\texttt{ContractingWith.fixedPoint}) for Theorem~\ref{thm:stoch_contracting}. A bounded
continuous $w$ solving the equation has $Tw = w$ pointwise (greatest elements are unique), so $w = V$ by uniqueness of the
fixed point. The rate is Mathlib's a priori estimate.
\end{proof}

\begin{theorem}[Optimal policy]
\label{thm:stoch_policy}
\lean{Econ.StochDynProg.policy, Econ.StochDynProg.policy_nonempty_compact_uhc}
\leanok
\uses{thm:stoch_bellman, lem:argmax_props, thm:uhc_argmax}
The policy $G(x, z) = \operatorname{argmax}_{y \in \Gamma(x,z)}[F + \beta E[V \mid z]]$ has nonempty compact values and
is upper hemicontinuous.
\end{theorem}
\begin{proof}
\leanok
Lemma~\ref{lem:argmax_props} and Theorem~\ref{thm:uhc_argmax} applied to the continuous objective of
Lemma~\ref{lem:stoch_objective}.
\end{proof}

\begin{remark}[Non-vacuity]
\label{rem:stoch_example}
\lean{Econ.StochExample.Q, Econ.StochExample.isStochastic_Q, Econ.StochExample.P}
\leanok
\uses{thm:stoch_bellman, thm:doeblin}
$X = [0, 1]$, two shock states switching with probability $1/2$, $\Gamma \equiv X$, $F = x z - y^2$ (bounded by $2$),
$\beta = 1/2$ satisfies every assumption, so Theorem~\ref{thm:stoch_bellman} applies. The fair-coin chain satisfies
Doeblin with $\varepsilon = 1/2$, so Theorem~\ref{thm:doeblin} gives a unique invariant distribution.
\end{remark}
